\section{Introduction}

A pseudoeffective divisor need not have a section in any positive
multiple.  This paper studies an analytic condition that makes such a
divisor effective up to rational linear equivalence, after a correction
on a prescribed boundary.  The condition is integrability of the
canonical divisor section for a positively curved anticanonical metric.
Its usefulness is birational: a correction on the base of a fibration
can disappear after pullback and pushforward to the original variety.

If the squared norm in a local frame of a singular metric $h$ is
$e^{-\phi}$, its multiplier ideal $\mathcal J(h)$ consists of the
holomorphic germs $u$ with $|u|^2e^{-\phi}$ locally integrable.
Thus $\mathcal J(h)\supset\mathcal O_Y(-C)$ means precisely that
the canonical section $s_C$ of the effective integral divisor $C$
has locally finite squared norm.  We call a rational divisor
\emph{rationally effective} if it is rationally linearly equivalent
to an effective rational divisor.  A rational divisor is
\emph{pseudoeffective} when its numerical class lies in the closure
of the cone generated by effective divisors.  On a smooth projective
variety this is equivalent to the existence of a singular Hermitian
metric with semipositive curvature current on its rational line bundle.

\begin{theorem}[Effectivity modulo a controlled boundary]
\label{thm:controlled-boundary}
Let $Y$ be a smooth projective complex variety of positive dimension
and let $C\ge0$ be an integral divisor.  Suppose $-K_Y+C$ has a
singular Hermitian metric $h$ whose curvature dominates a K\"ahler
form and for which
\[
 \mathcal J(h)\supset\mathcal O_Y(-C).
\]
For every pseudoeffective rational divisor $M$ on $Y$, there is an
effective rational divisor $C_0$ supported on $C$ such that
$M+C_0$ is rationally effective.
\end{theorem}

The support conclusion distinguishes the theorem from numerical
nonvanishing.  It permits the correction to be removed under a
specified birational map.  The analytic assumption also distinguishes
two notions of integrability: $s_C$ may have finite norm even when a
nonvanishing local frame does not.  The stronger, unweighted condition
will yield a separate finite-generation result.

\subsection{The boundary construction and its antecedents}

An effective rational boundary $\Delta$ is \emph{klt}
(Kawamata log terminal) if, on a log resolution $\sigma:Y'\to Y$,
every coefficient of
$\sigma^*(K_Y+\Delta)-K_{Y'}$ is less than one.
For a simple normal crossing boundary on a smooth variety, this
amounts to each boundary coefficient being less than one.

The proof connects multiplier-ideal approximation with big-boundary
nonvanishing.  Nadel vanishing~\cite{Nadel1990}, followed by very ample
regularity, makes the relevant multiplier-ideal systems globally
generated.  Point extension of Ohsawa--Takegoshi~\cite{OhsawaTakegoshi1987},
in the weighted Bergman approximation developed by
Demailly~\cite{Demailly2015}, compares the fixed divisor of a complete
system with the metric's singularities.  On a resolution
$\sigma:Y'\to Y$, weighted integrability produces a signed divisor
$\Psi=\Psi^+-\Psi^-$ whose coefficients are strictly below one.
Its negative part lies only over $C$ and the exceptional locus.

The approximation leaves an ample rational class $A$ in reserve.
For a prescribed pseudoeffective $M$, choose a small rational $t>0$
so that $A+tM$ remains ample.  General free and ample representatives
then turn the positive part of the signed divisor into an
\emph{effective big klt} boundary $\Delta$, with
\[
 K_{Y'}+\Delta\sim_{\mathbb Q}t\sigma^*M+\Psi^-.
\]
This is the pair to which the nonvanishing theorem of
Birkar--Cascini--Hacon--McKernan~\cite[Theorem D]{BCHM2010} applies.
A finite rational linear system upgrades its real-linear conclusion
to rational linear equivalence.  Pushforward gives the correction
$\sigma_*\Psi^-/t$, supported on $C$.
Sections~\ref{sec:approximation} and~\ref{sec:effectivity} give this
construction.  The latter also retains an analytic perturbation
argument: strong openness of Guan--Zhou~\cite{GuanZhou2015} and
small-exponent integrability of Skoda~\cite{Skoda1972} permit a small
pseudoeffective perturbation of the metric while preserving the
integrability of the distinguished section.

A different consequence begins with an unweighted integrable
semipositive metric that is smooth and strictly positive on one ordinary
open set.
Boucksom's volume criterion~\cite{Boucksom2002} first gives bigness;
a small mixture with a strict metric then gives a global curvature
lower bound while preserving integrability.  Section~\ref{sec:correction-ring}
makes this preparation explicit and permits perturbation in finitely
many arbitrary rational directions.  The resulting klt boundaries
fit the multigraded adjoint finite-generation results of
Cascini--Lazi\'c~\cite{CasciniLazic2012} and
Corti--Lazi\'c~\cite[Theorem 3.2(1)]{CortiLazic2013}.
More precisely, if $-K_Y+A_0$ has this unweighted metric, then for
every finite list of rational divisors $P_1,\ldots,P_s$ there are a
common rational $\epsilon>0$ and Cartier multiples
$B_i=k_0(A_0+\epsilon P_i)$ whose multigraded section ring is finitely
generated (Theorem~\ref{thm:correction-ring}).  This controls the
sections of specified Cartier divisors, not merely a numerical cone.

\subsection{Anticanonical nonvanishing and fixed-error conversion}

For a smooth projective variety $X$, anticanonical nonvanishing asks
whether some $H^0(X,-mK_X)$, $m>0$, is nonzero when $-K_X$ is nef.
Numerical effectivity is weaker: an effective divisor in the numerical
class need not supply a section of the actual line bundle.
For pairs $(X,\Delta)$, the corresponding divisor is
$-(K_X+\Delta)$, the anti-log-canonical divisor.
Lazi\'c--Matsumura--Peternell--Tsakanikas--Xie
\cite[Theorem~A]{LMPTX2023} established numerical nonvanishing for
projective log canonical threefold pairs with nef anti-log-canonical
divisor, assuming either that the variety is $\mathbb Q$-factorial
or that it has rational singularities.  M\"uller
\cite[Corollary~B]{Muller2025} proved nonvanishing for projective klt
threefold pairs with nef anti-log-canonical divisor.  His higher-dimensional
theorem assumes, in addition to nefness, that the anti-log-canonical
divisor of the induced pair on the general fiber of the maximal
rationally connected fibration is semiample \cite[Theorem~A]{Muller2025}.

Our principal conversion result uses smooth metric semipositivity and
works in every dimension.  Let $X$ be smooth connected projective,
let $L=-K_X$ carry a smooth semipositive Hermitian metric, and let $P$
be any fixed pseudoeffective Cartier divisor.  Then
Theorem~\ref{conv:fixed-error} gives the implication
\[
 H^0(X,m_jL-P)\ne0\ \text{for unbounded positive }m_j
 \quad\Longrightarrow\quad
 H^0(X,mL)\ne0\ \text{for some }m>0.
\]
The fixed error $P$ may itself have no effective representative.
A single Iitaka base converts the effective divisors into a
pseudoeffective rational divisor on that base.  The controlled-boundary
theorem makes this divisor rationally effective with correction on
base primes whose pullbacks disappear over $X$.  A rational function
from the base then corrects the original rational section without
introducing a pole on $X$.

Twisted differential forms supply the sequence through the determinant
method of Lazi\'c--Peternell \cite[Lemma~4.1]{LP2018}, adapted in
\cite[Lemma~5.1]{LMPTX2023}.  Ou's generic tangent-nefness theorem
\cite[Theorem~1.4]{Ou2023}, in the cotangent-subsheaf form
\cite[Theorem~4.1]{LMPTX2023}, makes the negative determinant
pseudoeffective.  Our extraction argument starts with positive
unbounded exponents and also includes numerically trivial $L$.
Thus nonzero sections of $(\Omega_X^1)^{\otimes p}\otimes m_jL$
for one fixed $p\ge0$ and unbounded positive $m_j$ give an
anticanonical section by Corollary~\ref{conv:tensor-conversion}.
The same holds for exterior forms.  This is an ordinary conversion;
invariant sections require further control of characters.

\subsection{Volume metrics and the distinct geometric refinements}

Let $X$ be smooth projective and let $L=-K_X$ carry a smooth
semipositive metric.  In a resolved rational fibration
\[
 X\xleftarrow{\ \pi\ } Z\xrightarrow{\ g\ }Y,
\]
assume $Z,Y$ are smooth projective, $\pi$ is birational, and $g$ is
surjective.  Write $K_{Z/X}=K_Z-\pi^*K_X$ for the relative canonical
(Jacobian) divisor.  Its canonical section belongs to $K_Z+\pi^*L$.
When
$\operatorname{rank}g_*\mathcal O_Z(K_{Z/X})=1$, its squared fiberwise
$L^2$ norm is a volume density $\rho$ on the base.  On the smooth
base-change open where this section generates the adjoint line,
$-\log\rho$ is a metric on $-K_Y$.
Berndtsson's direct-image theory supplies its positivity
\cite{Berndtsson2009,Berndtsson2011}.  Extending it over the boundary
with enough integrability is the geometric task preceding effectivity.

The main conversion proof uses a singular coefficient metric whose
curvature dominates a positive base form and whose total-space
multiplier ideal is trivial.  Singular adjoint direct-image positivity
of P\u aun--Takayama~\cite{PT2018} then gives a metric with a global
strict lower bound on $-K_Y+C$.  Its integral boundary has exceptional
pullback, and the distinguished section is locally square integrable.

Keeping the coefficient metric smooth gives a different output.
The coarea formula and a primitive-representative curvature calculation
produce an integral correction $C$ with $\pi_*g^*C=0$ and
\[
 |s_C|^2e^{-\phi}\in L^{1+\eta}_{\mathrm{loc}}
 \quad\text{for some }\eta>0.
\]
The metric is smooth and strictly positive on one ordinary open set.
The higher integrability exponent permits a mixture with a globally
strict metric while keeping $s_C$ integrable.  Thus the global bound
needed for Theorem~\ref{thm:controlled-boundary} follows from an
explicit additional step.

The companion \emph{Metric descent and rank-preserving contractions}
supplies the toroidal-volume construction used here
\cite[Theorem~2.1]{CompanionG}.  For a normal
projective model $p:W\to X$ and a surjective equidimensional toroidal
map $f:W\to Y$, its theorem assumes adjoint rank one at degree zero
on a smooth resolution and a patch, away from the Jacobian divisor,
where the pulled-back curvature dominates a positive base form.
It produces an effective rational correction $A_0$ with
\[
 p_*f^*A_0=0
\]
and a semipositive metric on $-K_Y+A_0$ with unweighted locally
integrable squared frame norms, smooth and strictly positive on a
nonempty ordinary open set.  That companion proves the toroidal
Jacobian calculation, the all-valuation integrability and the
boundary extension.  We use its metric after checking the stated
model, rank and curvature premises in each application.
The rational correction is used directly in
Corollary~\ref{conv:unweighted-effectivity}; no rounding or replacement
of its unweighted metric by an integral-boundary hypothesis is needed.

For the invariant problem, M\"uller \cite[Theorem~C]{Muller2025}
obtains naturally invariant sections for projective sub-log-canonical
pairs with semiample anti-log-canonical divisor and a commutative linear
algebraic group action.  Theorem~\ref{inv:nonvanishing} here assumes
$X$ smooth connected projective, $-K_X$ smoothly semipositive and
$\chi(X,\mathcal O_X)\ne0$.  For any algebraic torus acting on $X$,
it gives
\[
 H^0(X,-mK_X)^T\ne0\quad\text{for some }m>0
\]
with the natural tangent-determinant linearization.
The companion invariant-index theorem, including its actual value at
zero, supplies invariant cohomology
\cite[Theorem~1.1]{CompanionI}; compact-averaged hard Lefschetz
then gives forms with the coefficient shift $m+1$.
The local determinant and character arguments supply the sections
needed here.  The finite-cover structure and norm used in the
isometric application are those of the companion headline paper,
\cite[Theorem~2.1 and Lemma~2.2]{CompanionH}, which develops the geometric work of Demailly--Peternell--Schneider
and Campana--Demailly--Peternell~\cite{DPS1996,CDP2015}.
Our algebraic boundary criterion, smooth coarea proof and correction
ring are independent of that application.

Finally, a suitable auxiliary base has an effective big klt boundary
$\Delta$ with $K_Y+\Delta\sim_{\mathbb Q}C$ for an effective rational
$C$ supported on divisors disappearing over $X$.  When $\kappa(Y,C)=0$, a log
terminal model of the base is of Fano type: it admits an effective
rational boundary $\Theta$ with klt singularities such that
$-(K+\Theta)$ is ample.  This is a statement about
the auxiliary model.  Its proof retains the precise contracted-prime
support needed to return the resulting divisors to $X$.

The main route is the boundary and ring results followed by the
single-Iitaka conversion in Section~\ref{conv:section}.
Sections~\ref{coa:analytic-tools} and~\ref{jet:iitaka-section}
retain the smooth coarea estimate and direct jet-separation proof
of bigness; their applications reuse the common return of sections.
Section~\ref{fano:section} then develops the supported Fano-type
model.  The invariant base constructions in Sections~\ref{inv:section}
and~\ref{inv:dominated-alternative} lead to compact-monodromy descent
in Section~\ref{inv:isometric-descent}.
Finally, Section~\ref{ext:section} constructs invariant tensors from
extremal characters of the full equivariant Euler characteristic.
That construction is independent of the weight-zero index input and
descends tensors before applying ordinary conversion.
